\documentclass{article}

\usepackage[margin = 1in]{geometry}

\title{Ch. 30 CW \#1}
\author{Principles of Macroeconomics}
\date{\today}

\begin{document}

\maketitle

\begin{enumerate}

\item Isabella takes \$100 of currency from her wallet and deposits it into her checking account. If the bank adds the entire \$100 to reserves, the money supply \underline{\phantom{xxxxxxxxxx}}, but if the bank lends out some of the \$100, the money supply \underline{\phantom{xxxxxxxxxx}}.

	\begin{enumerate}
	
	\item increases; increases even more
	
	\item increases; increases by less
	
	\item is unchanged; increases
	
	\item decreases; decreases by less
	
	\end{enumerate}
	
\item If the reserve ratio is $\frac{1}{4}$ and the central bank increases the quantity of reserves in the banking system by \$120, the money supply increases by

	\begin{enumerate}
	
	\item \$90.
	
	\item \$150.
	
	\item \$160.
	
	\item \$480.
	
	\end{enumerate}

\item Bill Murray deposits \$500 that was in his sock drawer into a checking account at the local bank

	\begin{enumerate}
	
	\item How does the deposit initially (i.e. before any loans are made) change the T-account of the local bank? How does it change the money supply?
	
	\item If banks maintain a reserve ratio of 10\%, how will the local bank respond to the new deposit? 
	
	\item If every time the bank made a loan, the loan resulted in a new bank deposit equal to the amount of the loan, how much would the money supply in the economy expand in total?
	
	\item If banks instead maintained a reserve ratio of 5\%, how much would the money supply expand? 
	
	\end{enumerate}
	
\item Dan Aykroyd withdraws \$400 from his checking at the local bank and keeps it in his wallet.

	\begin{enumerate}
	
	\item If banks maintain a reserve ratio of 10\%, how will the withdrawal affect the money supply?
	
	\item If banks maintain a reserve ratio of 20\%, how will the withdrawal affect the money supply?
	
	\end{enumerate}

	\end{enumerate}

\end{document}